\documentclass[11pt]{article}
\usepackage[margin=1in]{geometry}
\usepackage{amsmath,amssymb,amsthm}
\usepackage{graphicx}
\usepackage{hyperref}
\usepackage{microtype}

\newtheorem{theorem}{Theorem}
\newtheorem{lemma}[theorem]{Lemma}
\newcommand{\ellone}{\ell_1}
\newcommand{\R}{\mathbb R}

\title{The exact maximal-equilateral number of\\uncountable $\ell_1(\Gamma)$}
\author{Candidate partial-result packet for arXiv:1109.5063}
\date{21 July 2026}

\begin{document}
\maketitle

\noindent\textbf{Status.} Candidate partial result, likely valid; human review requested.

\section{Source problem and exact gap}

For a normed space $X$, Swanepoel and Villa define $m(X)$ to be the minimum cardinality of a maximal equilateral subset of $X$ \cite{SV}. Their Conjecture~8 asks whether
\[
 m(X)\leq d+1
\]
for every $d$-dimensional normed space $X$. On page~4 they prove $m(\ellone^d)=4$ for $d\geq3$, explain why the smooth-point argument fails for uncountable $\ellone(\Gamma)$, and retain only
\[
 4\leq m(\ellone(\Gamma))\leq5
 \qquad(\Gamma\text{ uncountable}).
\]
The upper bound is their Theorem~6 (more explicitly, the $p=1$ construction in Proposition~19). The result below closes this exact-value gap. It is adjacent progress rather than a solution of the finite-dimensional conjecture.

\begin{center}
\includegraphics[width=\textwidth]{figures/open_problem_crop.png}
\end{center}

\section{Proof intuition}

Four points in an $\ell_1$ space admit a useful coordinatewise cut decomposition. Sort the four scalar values in each coordinate. The bottom and top gaps act as singleton cuts, while the middle gap acts as one of the three possible $2$--$2$ cuts. Equilateral distances force all four total singleton weights to agree and all three total $2$--$2$ weights to agree.

Now take, in every coordinate, the midpoint of the two middle scalar values. Its distance to a given point is one half of every middle gap, plus precisely the outer gaps assigned to that point. The cut equalities therefore make this a common center. Its radius is strictly smaller than the equilateral side. Four $\ell_1(\Gamma)$ vectors use only countably many coordinates, so when $\Gamma$ is uncountable an unused coordinate can supply the missing radius and extend the four-point set.

For the reverse inequality, the five-point Hadamard configuration used in the source at $p=1$ is shown directly to remain maximal even when arbitrarily many new coordinates are available.

\section{The four-point center lemma}

\begin{lemma}\label{lem:center}
Let $a_1,a_2,a_3,a_4$ be a $\lambda$-equilateral subset of $\ellone(\Gamma)$. There exist $c\in\ellone(\Gamma)$ and $r<\lambda$ such that
\[
 \|a_i-c\|_1=r\qquad(i=1,2,3,4).
\]
\end{lemma}

\begin{proof}
Fix a coordinate $\gamma$. Sort the four real numbers $a_i(\gamma)$ as
\[
 b_1\leq b_2\leq b_3\leq b_4
\]
and write their consecutive gaps as
\[
 \alpha_\gamma=b_2-b_1,\qquad
 \beta_\gamma=b_3-b_2,\qquad
 \delta_\gamma=b_4-b_3.
\]
If $\alpha_\gamma>0$, exactly one label attains $b_1$; assign this singleton gap to that label. If $\delta_\gamma>0$, assign it to the unique label attaining $b_4$. If $\beta_\gamma>0$, the lower two and upper two labels define one of the three unordered partitions
\[
 12|34,\qquad13|24,\qquad14|23.
\]
Zero gaps may be assigned arbitrarily. All gap series below converge, since each coordinate range is at most $\sum_i|a_i(\gamma)|$.

Let $s_i$ be the total outer-gap weight assigned to label $i$, and let $q_{12},q_{13},q_{14}$ be the total middle-gap weights of the three displayed partitions. A scalar difference crosses exactly those consecutive gaps whose cuts separate its two labels. Hence
\begin{align*}
 \|a_1-a_2\|_1&=s_1+s_2+q_{13}+q_{14},&
 \|a_3-a_4\|_1&=s_3+s_4+q_{13}+q_{14},\\
 \|a_1-a_3\|_1&=s_1+s_3+q_{12}+q_{14},&
 \|a_2-a_4\|_1&=s_2+s_4+q_{12}+q_{14},\\
 \|a_1-a_4\|_1&=s_1+s_4+q_{12}+q_{13},&
 \|a_2-a_3\|_1&=s_2+s_3+q_{12}+q_{13}.
\end{align*}
Equality of complementary pair distances first gives
\[
 s_1+s_2=s_3+s_4,\qquad
 s_1+s_3=s_2+s_4,\qquad
 s_1+s_4=s_2+s_3.
\]
Thus $s_1=s_2=s_3=s_4=:s$. Comparing the remaining pair distances then gives
\[
 q_{12}=q_{13}=q_{14}=:q.
\]
Consequently the common side length is
\[
 \lambda=2s+2q.\tag{1}
\]

Define
\[
 c(\gamma):=\frac{b_2+b_3}{2}.
\]
This defines an element of $\ellone(\Gamma)$ because
$|c(\gamma)|\leq\max_i|a_i(\gamma)|\leq\sum_i|a_i(\gamma)|$.
At coordinate $\gamma$, every label pays $\beta_\gamma/2$ to reach $c(\gamma)$; the unique bottom label additionally pays $\alpha_\gamma$, and the unique top label additionally pays $\delta_\gamma$. Therefore
\[
 \|a_i-c\|_1=s_i+\frac12(q_{12}+q_{13}+q_{14})
 =s+\frac32q=:r
\]
for every $i$. By (1),
\[
 \lambda-r=s+\frac12q>0,
\]
where strictness follows from $\lambda>0$ and $s,q\geq0$.
\end{proof}

\section{Exact value for uncountable \texorpdfstring{$\ell_1(\Gamma)$}{ell1(Gamma)}}

\begin{theorem}\label{thm:main}
For every uncountable set $\Gamma$,
\[
 \boxed{m(\ellone(\Gamma))=5.}
\]
\end{theorem}

\begin{proof}
First consider a four-point $\lambda$-equilateral set $A=\{a_1,a_2,a_3,a_4\}$. By Lemma~\ref{lem:center}, choose $c$ and $r<\lambda$ equidistant from its points. The union of the supports of $a_1,\ldots,a_4$ is countable, and the construction of $c$ uses no other coordinates. Choose $\gamma_0\in\Gamma$ outside this union and put
\[
 x=c+(\lambda-r)e_{\gamma_0}.
\]
The new coordinate is disjoint from every $a_i-c$, so
\[
 \|x-a_i\|_1=\|c-a_i\|_1+\lambda-r=\lambda
\]
for all $i$. Thus no four-point equilateral set is maximal. The source recalls Petty's theorem that in dimension at least three every equilateral triple extends to four points \cite{SV}; pairs extend to triples by the standard continuity argument. Hence $m(\ellone(\Gamma))\geq5$.

For the matching upper bound, choose four distinct coordinates of $\Gamma$ and, suppressing all remaining zero coordinates, set
\begin{align*}
 p_1&=(1,1,1,0),&p_2&=(1,-1,-1,0),\\
 p_3&=(-1,1,-1,0),&p_4&=(-1,-1,1,0),\\
 p_5&=(0,0,0,1).
\end{align*}
All ten pairwise $\ell_1$ distances are $4$.

Suppose that $x=(x_1,x_2,x_3,x_4,z)$ were at distance $4$ from all five points, where $\rho=\|z\|_1$ accounts for all remaining coordinates. Put
\[
 h(t)=|t-1|-|t+1|.
\]
Equality of the distances from $x$ to $p_1,p_2,p_3,p_4$ gives
\[
 h(x_2)+h(x_3)=0,\qquad
 h(x_1)+h(x_3)=0,\qquad
 h(x_1)+h(x_2)=0.
\]
These equations force $h(x_1)=h(x_2)=h(x_3)=0$, and $h(t)=0$ only for $t=0$. Hence $x_1=x_2=x_3=0$. The distance to any of the first four points now yields
\[
 3+|x_4|+\rho=4,\qquad\text{so }|x_4|+\rho=1.
\]
But then
\[
 \|x-p_5\|_1=|x_4-1|+\rho
 \leq |x_4|+1+\rho=2<4,
\]
a contradiction. Thus the five-point set is maximal and $m(\ellone(\Gamma))\leq5$.
\end{proof}

\section{Verification, scope, and novelty}

The script \texttt{code/verify.py} solves the cut-weight linear system symbolically. It confirms
\[
 s_1=s_2=s_3=s_4,\qquad q_{12}=q_{13}=q_{14},\qquad
 \lambda-r=s+\frac q2.
\]
It also checks exact rational permutation-symmetrized examples with three nonzero gaps per coordinate, the unused-coordinate extension, and all ten distances in the maximal five-point set.

This theorem determines the exact value left between the source's lower and upper observations for uncountable $\ellone(\Gamma)$. It does not resolve the finite-dimensional Conjecture~8.

Cheap run indexes were searched for arXiv:1109.5063 and the core maximal-equilateral terminology. Bounded web searches on 21 July 2026 used exact variants of ``$m(\ell_1(\Gamma))=5$'', ``uncountable $\ell_1$ maximal equilateral'', and four-point $L_1$ cosphericity. They found the source, related infinite-dimensional work, and an expository five-point noncospherical example, but no later proof of this exact value or the four-point lemma. This is not an exhaustive novelty certification.

\medskip
\noindent\textbf{Human-review recommendation.} Check the cut bookkeeping at tied scalar values, convergence of the nonnegative gap sums, and the strict radius calculation. The five-point maximality argument is independent and elementary.

\begin{thebibliography}{1}
\bibitem{SV}
K.~J. Swanepoel and R. Villa,
\emph{Maximal equilateral sets},
Discrete Comput. Geom. \textbf{50} (2013), 354--373;
arXiv:\href{https://arxiv.org/abs/1109.5063}{1109.5063}.
\end{thebibliography}

\end{document}
